\documentclass{article}

\title{Ch. 14 HW \#2}
\author{Principles of Microeconomics}
\date{\today}

\begin{document}

\maketitle

\subsection*{Instructions} This is not the entirety of Ch. 14 HW \#2. This is just the portion of the assignment that we started in class.

\subsection*{Problems}

\begin{enumerate}

\item Buffy is thinking about opening an amulet store. She esstimates that it would cost \$350,000 per year to rent the location and buy the merchandise. In addition, she would have to quit her \$80,000 per year job as a vampire hunter.

	\begin{enumerate}
	
	\item Define \textbf{opportunity cost}.
	
	\item What is Buffy's opportunity cost of running the store for a year?
	
	\item Buffy thinks she can sell \$400,000 worth of amulets in a year. What would her accountant consider the store's profit?
	
	\item Should Buffy open the store? Explain.
	
	\item How much revenue would the store need to generate for Buffy to earn positive economic profit?
	
	\end{enumerate}
	
\item A commercial fisherman notices the following relationship between hours spent fishing and the quantity of fish caught:

	\begin{center}
	\begin{tabular}{ c | c }
	\textbf{Hours} & \textbf{Quantity of Fish (in pounds)} \\
	\hline
	0 & 0 \\
	1 & 10 \\
	2 & 18 \\
	3 & 24 \\
	4 & 28 \\
	5 & 30 \\
	\end{tabular}
	\end{center}
	
	\begin{enumerate}
	
	\item What is the marginal product of each hour spent fishing?
	
	\item Use these data to graph the fisherman's production function. Explain its shape.
	
	\item The fisherman has a fixed cost of \$10 (his pole). The opportunity cost of his time is \$5 per hour. Graph the fisherman's total-cost curve. Explain its shape.
	
	\end{enumerate}

\end{enumerate}

\end{document}